\chapter*{Conclusion} % chapter* je necislovana kapitola
\addcontentsline{toc}{chapter}{Conclusion} % rucne pridanie do obsahu
\markboth{Conclusion}{Conclusion} % vyriesenie hlaviciek

This thesis introduced \emph{Global Assignment TETRIS} (GA-TETRIS), a permutation search algorithm for block pruning of transformer weight matrices. The central idea is that finding the best permutation at a fixed block mask is exactly the linear assignment problem and can be solved in polynomial time, where the Original TETRIS uses a greedy single-swap heuristic that stops at the first local optimum. To this exact step we added two heuristics --- multiplicative noise injection during the assignment phase and a second phase of random-swap-with-refinement --- to escape local optima created by the iterative re-masking. We also introduced two simple heuristics, Sort-by-Norm and Random-Swaps, to compare against the more complex TETRIS algorithms.

We compared GA-TETRIS against the Original TETRIS, Sort-by-Norm, Random-Swaps baselines, and the no-permutation Block-Wanda baseline on two Vision Transformers (0.9M and 13.4M parameters) and SmolLM2-360M. The per-layer evaluation showed that GA-TETRIS produces the largest score improvements at wide block widths consistently across all three models, with the gains concentrated in the early transformer blocks. The crossover block width at which GA-TETRIS overtakes simple Sort-by-Norm was approximately the same across models, suggesting it is a property of the Wanda-score distribution rather than of model size.

The whole-model picture was less favorable. Performance collapsed sharply with block width on all three models, with only the Larger ViT at block $1 \times 2$ preserving substantial accuracy. In that single regime GA-TETRIS achieved the highest accuracy, beating the next-best algorithm by $1.85$ percentage points and the no-permutation baseline by $2.6$ percentage points. On SmolLM2 at the same block width, however, surprisingly the simple Sort-by-Norm produced the lowest perplexity, with both Original TETRIS and GA-TETRIS noticeably worse. The compute cost of GA-TETRIS is also substantial: at block $1 \times 2$ on the Larger ViT it ran approximately $7.6$ seconds per prunable layer, against essentially zero for Sort-by-Norm. Whether the additional accuracy is worth the compute depends on the deployment context.

As future work, several directions remain open. We did not fine-tune the pruned models, which is standard practice in the pruning literature and might substantially recover task quality. We also operated at reduced compute budgets in places: GA-TETRIS used $T_{phase2} = 100$, and Random-Swaps used $T_{swaps} = 1000$. We saw a good evidence that the algorithms didn't converge at the given random swap limits and could potentially benefit from more iterations at the cost of multiple hours compute time on larger models. Apart from this, the sharp gap between per-layer score improvement and whole-model task quality, especially on SmolLM2, suggests that the entrywise $L_1$ pruning objective is not perfectly aligned with what determines downstream behavior. Finally, the algorithm has hyperparameters we did not tune systematically, in particular the swap fraction $\rho$ and the initial noise scale $\eta_0$.
